\chapter{Introduction} \label{cha:introduction}
\pagestyle{plain}

In the age of widespread digital communication and increasing concerns about data privacy, 
the methods of concealing information have become vital in computer and systems sciences research. 

Steganography, derived from the Greek word \textit{steganos} (covered) and the word \textit{graphein} (writing), 
is defined as the practice of concealing information inside often-innocent digital media, such as images, 
for secret communication \cite{Fridrich2002}.
While cryptography encrypts information in order to make it unreadable, steganography conceals the 
very existence of the information \cite{Fridrich2002}, making it valuable for 
scenarios where the attention drawn to encrypted messages would be undesirable.

Steganographic methods have a number of serious implications concerning security and forensics. 
Although they can be used for legal purposes, such as digital watermarking used for copyright protection, 
secure communication in countries that ban or restrict encryption, these methods can be used for malicious tasks as well \cite{Fridrich2002, Eaton2016}.
Malicious actors have been known to use steganography to covertly communicate, conceal incriminating evidence, and 
distribute illicit content such as terrorist training manuals and child pornography \cite{Eaton2016}.
This makes the ability to detect steganographically hidden data important to law enforcement, 
counterintelligence, and other cybersecurity practitioners and digital forensic investigators \cite{Eaton2016}.

The countermeasure to steganography is steganalysis.
Its objective is to determine whether digital media contains embedded data \cite{Fridrich2002}.
Steganography and steganalysis are adversarial fields in which developments in one drive developments in the other.
Although steganalysis methods can be classified in several different ways \cite{Selvaraj2021}, 
this study focuses on two different approaches to steganalysis: traditional statistical techniques and deep-learning-based techniques.

Traditional statistical methods, such as Regular-Singular (RS) analysis, detect disruptions in the spatial correlations and 
local smoothness that LSB embedding introduces into pixel data \cite{Fridrich2001, Manoharan2008}.

Deep-learning-based steganalysis uses machine learning models, such as Convolutional Neural Networks (CNNs), 
to detect the subtle, high-dimensional patterns and noise artifacts introduced by steganographic embedding \cite{Qian2015}. 

The advancements in deep learning research have resulted in a paradigm shift for steganalysis research with a 
tendency to regard Neural Networks as the norm for solving steganalytic detection tasks \cite{Qian2015, Selvaraj2021}. 
However, this shift might neglect computational complexity and the baseline detection accuracies of older methods.
Although the use of Convolutional Neural Networks (CNNs) with state-of-the-art performance for complex 
and adaptive steganographic attacks is compelling, their superiority over traditional methods for basic LSB embedding schemes 
within lossless images is still worthy of research.

This study compares the two approaches, and aims to provide a direct and empirical comparison of their respective effectiveness.


\section{Delimitations}

\subsection{Image file format (PNG Only)}
The study focuses on Least Significant Bit (LSB) steganography, which requires exact pixel value preservation. 
Lossy image formats, such as JPEG, destroy LSB information in the compression process and introduce artifacts that interfere with the 
specific statistical analysis methods being tested. Performing both steganography and steganalysis on such image formats requires significantly higher sophistication,
and these are outside our scope.

\subsection{Spatial Domain LSB Substitution}
This research aims to compare the baseline effectiveness of a statistical method (RS-analysis) against a Deep Learning model (CNN) in the spatial domain. 
While complex, highly adaptive steganographic algorithms (e.g., HUGO \cite{pevny2010using}, S-UNIWARD \cite{holub2014universal}) fall outside the scope of this comparison due to 
requiring entirely different steganalytic approaches, foundational spatial techniques, such as basic Sobel-based edge-adapted steganography, is included to provide a comprehensive baseline.

\subsection{Dataset and Color Channels}
The BOSSBase dataset is widely used in the field of steganalysis research, ensuring reproducibility \cite{bas2011break, Sedighi2016}. 
The dataset's images are in a lossless format (.png), and as the images in the dataset are in grayscale, 
the embedding artifacts are isolated within the luminance values.

\subsection{Payload Granularity}
This study opts for discrete payload sizes rather than a continuous spectrum. 
In order to answer the sub-question regarding changes as payload decreases, 
the study tests fixed embedding rates (25\%, 50\%, and 75\% of image embedding capacity for respective steganographic method) 
to provide clear comparative bins for the True Positive Rate and overall detection accuracy analysis.


\section{Research Question}

This study is an empirical comparison, intent on investigating the effectiveness of traditional statistical steganalysis versus modern 
deep learning classification in detecting steganographic embedding. 

This objective is condensed into the following \textbf{Research Question (RQ)}:

\begin{itemize}
  \item[] \textbf{RQ:} How does Convolutional Neural Network (CNN) compare to a RS-Analysis statistical test in accurately and reliably detecting LSB steganography in PNG images?
\end{itemize}

With the help of the following \textbf{Sub-Question (SQ)}:

\begin{itemize}
  \item[] \textbf{SQ :} How does the detection accuracy of each method change as the size of the hidden payload decreases?
\end{itemize}
